\documentclass[11pt]{article}
\usepackage{coursenotes}
\begin{document}
\notesheader{8}{Double Machine Learning}{Flexible controls without biased answers}

\begin{decision}
A price, fee, or discount was set by a system reacting to demand signals. The firm wants the effect of the price itself: an
elasticity to hold against a break-even. Many signals must be controlled for, in complicated ways, so machine learning is
needed. But a model that predicts outcomes well does not, by itself, give an unbiased effect. These notes show how to use
machine learning for the controls without letting its errors bias the causal estimate, and how to say how wrong the controls would
have to be to change the decision.
\end{decision}

\begin{goals}
  \item state the partially linear model, its identifying assumptions, and its estimand when effects vary;
  \item derive the residual-on-residual (Frisch--Waugh--Lovell) form and explain why both sides must be residualized;
  \item explain the two failures of plug-in machine learning and how orthogonality and cross-fitting fix them;
  \item compute a DML estimate and its (clustered) standard error, and diagnose overlap with a continuous treatment;
  \item use the interactive model for an ATE, and run an omitted-variable sensitivity analysis with benchmarking.
\end{goals}

%=====================================================================================
\sectionone

\subsection{The partially linear model}

With outcome $Y$, treatment $D$ (binary or continuous, such as a log price), and controls $X$,
\[
  Y = \theta D + g(X) + \varepsilon, \qquad \E[\varepsilon \mid D, X] = 0 .
\]
The model is linear in $D$, so $\theta$ is one number (an elasticity if $Y$ and $D$ are logs), and unrestricted in $X$: $g$ may be
any function. Write the \emph{nuisance functions} $\ell(X) = \E[Y \mid X]$ and $r(X) = \E[D \mid X]$ (the propensity score when $D$ is
binary).

\begin{assum}{Conditional mean independence}{cmi}
$\E[\varepsilon \mid D, X] = 0$: given $X$, nothing left in the outcome moves with $D$.
\end{assum}
\begin{assum}{Residual treatment variation}{rtv}
$\E[\Var(D \mid X)] > 0$: treatment is not a deterministic function of $X$.
\end{assum}

The first is Meeting~7's exchangeability for a possibly continuous treatment. When a recorded rule set the price, it reads: the
rule's inputs are all in $X$, nothing unrecorded moved the price, and $X$ is measured \emph{as of the moment the price was set}. A
signal recorded later (booking pace after the price posted, a competitor index that reacts to our price) is partly caused by $D$:
a mediator or collider, and a bad control.

\begin{prop}{Frisch--Waugh--Lovell in population form}{fwl}
Under the model, $Y - \ell(X) = \theta\{D - r(X)\} + \varepsilon$, so
\[
  \theta = \frac{\E[(D - r(X))(Y - \ell(X))]}{\E[(D - r(X))^2]} .
\]
\end{prop}
\begin{proof}
Taking $\E[\,\cdot \mid X]$ gives $\ell(X) = \theta r(X) + g(X)$; subtract from the model. Multiply by $D - r(X)$ and take expectations;
the $\varepsilon$ term vanishes.
\end{proof}

Note that $\ell \ne g$: predicting $Y$ from $X$ estimates $\ell = \theta r + g$, which already contains the part of the treatment $X$
explains. That is why \emph{both} sides are residualized. Subtracting $\ell$ from $Y$ and regressing on raw $D$ discards the very
treatment variation the estimate needs.

\begin{prop}{What $\theta$ estimates when effects vary}{vw}
If the true slope varies, $Y = \theta(X)D + g(X) + \varepsilon$, the residual-on-residual coefficient converges to
$\E[\theta(X)\Var(D \mid X)]/\E[\Var(D \mid X)]$.
\end{prop}
\begin{proof}
$\E[(D - r)(Y - \ell) \mid X] = \theta(X)\Var(D \mid X)$ (the $g$ and $\varepsilon$ terms vanish) and $\E[(D-r)^2 \mid X] = \Var(D \mid X)$;
average over $X$.
\end{proof}

$\theta$ is a weighted average, weighted by how much treatment still varies once $X$ is known. Units the system priced almost
deterministically contribute almost nothing. For binary $D$ the weights are $e(X)(1-e(X))$: this is not the ATE (Meeting~7).
Report which weighted average you have, and for whom.

\subsection{Why plugging in machine learning fails}

\paragraph{Regularization bias.} Learners trade bias for variance. A covariate that strongly predicts $D$ but adds little to
predicting $Y$ may be shrunk out of an outcome model fitted on $(D, X)$, or under-used by a tree ensemble once $D$ is available.
Its confounding is then attributed to $D$. A better prediction of $Y$ is not a better estimate of $\theta$.

\paragraph{Overfitting bias.} If the nuisance models are fitted on the rows they later score, residuals are too small and their
errors are correlated through the shared rows. That correlation lands on $\theta$, grows with the learner's flexibility, and is
invisible from inside the fit.

\subsection{Cross-fitting and the product of errors}

\begin{prop}{The DML algorithm}{dml}
(1) Split rows into $K$ folds (often 5), at the level of any dependence (clusters, dates). (2) For each fold, fit $\hat\ell$ and $\hat r$
on the other folds and predict this fold, giving out-of-fold residuals $\tilde Y_i = Y_i - \hat\ell(X_i)$ and $\tilde D_i = D_i - \hat r(X_i)$.
(3) Estimate
\[
  \hat\theta = \frac{\sum_i\tilde D_i\tilde Y_i}{\sum_i\tilde D_i^2}, \qquad
  \widehat\SE(\hat\theta) = \frac{\sqrt{\sum_i\tilde D_i^2\hat\varepsilon_i^2}}{\sum_i\tilde D_i^2}, \quad \hat\varepsilon_i = \tilde Y_i - \hat\theta\tilde D_i .
\]
With clustered data replace $\sum_i\tilde D_i^2\hat\varepsilon_i^2$ by $\sum_c(\sum_{i \in c}\tilde D_i\hat\varepsilon_i)^2$.
\end{prop}
\begin{proof}
The estimator is the sample analog of Result~\ref{prop:fwl}; the SE is the heteroskedasticity-robust (or cluster-robust) SE of
the residual regression. Its validity rests on the bias calculation below and the rate conditions in Section~3.
\end{proof}

Why it works: write the nuisance errors $e_r = \hat r - r$, $e_\ell = \hat\ell - \ell$, and $\nu = D - r(X)$. Then
$\tilde D = \nu - e_r$ and $\tilde Y = \theta\nu + \varepsilon - e_\ell$. Cross-fitting makes the errors independent of the fold being
scored, so the cross terms with $\nu$ and $\varepsilon$ vanish in expectation, leaving
\[
  \hat\theta \approx \theta + \frac{\E[e_re_\ell] - \theta\,\E[e_r^2]}{\E[\nu^2] + \E[e_r^2]} .
\]
The first term is a \emph{product} of two errors: if each model is moderately wrong, the product is much smaller than either. This
is \emph{Neyman orthogonality}: first-order errors in either nuisance do not move the estimating equation. The second term, the
squared error of the treatment model, attenuates $\hat\theta$ toward zero and gets no product protection; it vanishes only if $\hat r$ is
accurate. Diagnose it with the share of $\sum\tilde D^2$ that is model error, and fix it with a better treatment model or more data.

Repeat the whole procedure over several random fold splits and report the median, adding the spread across splits to the SE. Choose
each nuisance learner by its own cross-validated prediction error (that is legitimate: they are prediction problems), and report at
least two. Disagreement means the answer depends on modeling choices.

\subsection{Overlap with a continuous treatment}

Wherever $\Var(D \mid X)$ is near zero (a pricing rule applied without deviation), there is no comparison, and those rows contribute
nothing however many there are. Two diagnostics: the out-of-fold $R^2$ of the treatment model (near 1 means little identifying
variation is left) and the share of $\sum\tilde D_i^2$ from each region (which says whom the estimate describes). A local slope supports
decisions near the prices the data contain; a formula that extrapolates it (such as an optimal-markup rule) can point far outside them.

\subsection{The interactive model and the ATE}

For binary $D$ and the ATE, use $Y = g(D, X) + \varepsilon$ and the cross-fitted AIPW score of Meeting~7:
\[
  \hat\theta_{\ATE} = \frac1n\sum_i\phi_i, \quad
  \phi_i = \hat\mu_1(X_i) - \hat\mu_0(X_i) + \frac{D_i(Y_i - \hat\mu_1(X_i))}{\hat e(X_i)} - \frac{(1-D_i)(Y_i - \hat\mu_0(X_i))}{1 - \hat e(X_i)},
\]
with standard error $\text{sd}(\phi_i)/\sqrt n$. The score is orthogonal in the same sense: its bias is a product of the outcome and
propensity models' errors. The choice between the partially linear and interactive models is a choice of estimand. Regressing $\phi_i$ on
a few modifiers gives the best linear projection of the CATE on them, with valid SEs; regressing it on $X$ flexibly is the DR-learner.

\subsection{Sensitivity to unmeasured confounding}

\begin{prop}{An omitted-variable bound}{ovb}
Suppose an omitted confounder $U$ would explain a share $R^2_{D}$ of the residual treatment variance (after $X$) and a share $R^2_{Y}$ of
the residual outcome variance (after $D$ and $X$). Then the bias in the partially linear $\hat\theta$ is at most
\[
  |\text{bias}| \le \sqrt{\frac{R^2_{Y}\,R^2_{D}}{1 - R^2_{D}}}\;\cdot\;\frac{\text{sd}(\tilde Y - \theta\tilde D)}{\text{sd}(\tilde D)} .
\]
\end{prop}
\begin{proof}
This is the omitted-variable-bias formula in partial-$R^2$ form (Cinelli and Hazlett, 2020) applied to the residualized regression; the
bound uses $|\text{corr}| \le 1$ for the part of $U$ not explained by $D$ and $X$.
\end{proof}

The two $R^2$ values are unknowable, so use the bound in two ways. The \emph{robustness value}: the equal $R^2$ at which the estimate
reaches the break-even. \emph{Benchmarking}: drop an observed control, re-estimate, and compute how much residual variance it explained;
then ask whether an omitted factor could be that strong, or a fraction of it. A conclusion that survives a confounder as strong as the
strongest observed control is robust; one that flips at a fraction of a minor control is not.

\begin{pitfall}[DML errors]
Reading a slope off a model of $Y$ on $(D, X)$; fitting nuisances on the rows they score; random folds on clustered data; controls
recorded after the price was set; reporting a variance-weighted $\theta$ as the ATE; extrapolating a local elasticity to prices never
charged.
\end{pitfall}

%=====================================================================================
\sectiontwo

\paragraph{Setting.} QuickBite's system sets each zone's delivery fee by the hour from demand and courier-supply signals. Should the
standard fee rise from ¥4.0 to ¥4.4? Margin per order excluding the fee is ¥6. Profit per zone-hour is $Q(f)(6 + f)$, whose elasticity
with respect to $f$ is $\varepsilon + f/(6 + f) = \varepsilon + 0.4$ at $f = 4$. A rise pays only if the order elasticity $\varepsilon$ exceeds $-0.4$.
The data are 40{,}000 zone-hours from a simulation calibrated to this setting (script \texttt{checks/m08\_sim.py}; true elasticity
$-0.6$), with $Y$ log orders, $D$ log fee, and $X$ hour, rain, weekend, zone, and supply signals.

\paragraph{Step 1: naive.} The regression of log orders on log fee gives $+1.20$: fees are high exactly when demand is high (dinner,
rain), so high-fee hours have many orders for reasons unrelated to the fee. Taken at face value it says to raise the fee.

\paragraph{Step 2: DML by hand on six rows.} Out-of-fold residuals:
\begin{center}
\begin{tabular}{lcccccc}
\toprule
Zone-hour & 1 & 2 & 3 & 4 & 5 & 6 \\
\midrule
$\tilde D$ & 0.12 & $-0.05$ & 0.08 & $-0.10$ & 0.02 & $-0.07$ \\
$\tilde Y$ & $-0.09$ & 0.01 & $-0.03$ & 0.08 & $-0.02$ & 0.03 \\
\bottomrule
\end{tabular}
\end{center}
$\sum\tilde D\tilde Y = -0.0242$ and $\sum\tilde D^2 = 0.0386$, so $\hat\theta = -0.63$ on these rows.

\paragraph{Step 3: full sample.} Five-fold DML with gradient boosting gives $\hat\theta = -0.575$ (SE $0.015$); with lasso nuisances, $-0.524$
(SE $0.015$). The fee model's out-of-fold $R^2$ is $0.74$, so a quarter of fee variation is left to identify the effect. Both estimates are
smaller in magnitude than the truth: the attenuation term $-\theta\E[e_r^2]$ at work, larger for the lasso, whose fee model is more misspecified. Both are
well below $-0.4$: at $-0.575$ a 10\% fee rise changes profit by a factor of about $(1.1)^{-0.575} \times 10.4/10 = 0.985$, a 1.5\% loss.

\begin{center}
\begin{tikzpicture}
\begin{axis}[width=0.62\linewidth, height=5.2cm, xlabel={$\tilde D$ (log-fee residual)}, ylabel={$\tilde Y$ (log-order residual)},
  axis lines=left, xmin=-0.35, xmax=0.35, ymin=-1.1, ymax=1.1, xticklabel style={/pgf/number format/fixed}]
\addplot[only marks, mark=*, mark size=1.1pt, color=noteblue] table[col sep=comma, x index=0, y index=1] {m08_resid_sample.csv};
\addplot[only marks, mark=square*, mark size=1.6pt, color=warnred] table[col sep=comma, x index=0, y index=1] {m08_resid_bins.csv};
\addplot[thick, sbgreen, domain=-0.35:0.35] {-0.575*x};
\end{axis}
\end{tikzpicture}

\small Figure 1. 300 randomly chosen out-of-fold residual pairs (dots), the means of all 40{,}000 pairs in 20 bins of $\tilde D$ (squares),
and the DML slope ($-0.575$). Single points are noisy; the binned means show the slope the estimate uses.
\end{center}

\paragraph{Step 4: sensitivity.} The residual SDs give a ratio $\text{sd}(\tilde Y - \hat\theta\tilde D)/\text{sd}(\tilde D) = 0.304/0.101 = 3.0$. A confounder with
$R^2 = 0.05$ on both sides could move the estimate by at most $\sqrt{0.05 \times 0.05/0.95} \times 3.0 = 0.15$, giving $[-0.73, -0.42]$: still below
$-0.4$. Reaching the break-even needs a bias of $0.175$, an equal $R^2$ of about $0.057$. Benchmark: rain explains $0.18$ of the residual fee
variance and $0.32$ of the residual order variance. An omitted factor a third as strong as rain on the fee side would overturn the
conclusion. Unrecorded local events could plausibly be that strong.

\paragraph{Recommendation.} Do not raise the fee. Because a moderately strong unrecorded confounder could overturn the conclusion,
confirm it with a randomized fee test (a switchback by hour in a few zones, which respects courier interference; Handout H1).

\begin{exercise}[computation]
Using the DML estimate $-0.575$, what fee change does the model suggest, and why should you not compute an ``optimal fee'' from the constant-elasticity formula?
\end{exercise}
\answer{Since $\varepsilon + f/(6+f) < 0$ at $f = 4$, profit rises if the fee falls slightly: the local evidence supports a small cut, not a rise.
The constant-elasticity optimum for a cost-plus structure would extrapolate the local slope to fees the system rarely charged (the
treatment model leaves little variation far from current fees), so the formula's answer is outside the support. The supported action is a
small, randomized test of a cut.}

\begin{exercise}[judgment]
A colleague suggests adding ``orders in the same zone in the previous 15 minutes'' and ``competitor fee in the same hour'' to $X$. Assess
each.
\end{exercise}
\answer{Previous-15-minute orders: fine only if measured before this hour's fee was set and not affected by it; if the window overlaps the
fee's posting it is partly caused by $D$ (a mediator). Competitor fee: if competitors react to QuickBite's fee within the hour, it is
post-treatment and a bad control; if it is set independently beforehand, it is a legitimate confounder. Check timestamps against the fee-setting
time.}

%=====================================================================================
\sectionthree

\begin{extension}[Orthogonality, rates, and Robinson]
A score is Neyman orthogonal if its derivative with respect to the nuisances is zero at the truth. Valid $\sqrt n$ inference needs the product of
nuisance error rates to be $o(n^{-1/2})$ (each $o(n^{-1/4})$ suffices) and the treatment model's squared error $o(n^{-1/2})$. Robinson (1988) proposed the
residual-on-residual estimator with kernel nuisances; cross-fitting removed the entropy conditions that restricted that theory, allowing
high-dimensional learners.
\end{extension}

\begin{extension}[Post-double-selection and dependent data]
The linear predecessor: lasso $Y$ on $X$, lasso $D$ on $X$, take the union of selected variables, and regress $Y$ on $D$ and the union (Belloni,
Chernozhukov, and Hansen, 2014). With panel or time-series data, form folds by cluster or forward in time, and check stability with coarser folds.
\end{extension}

\subsection*{Annotated reading}
\begin{readinglist}
\reading{Chernozhukov, Chetverikov, Demirer, Duflo, Hansen, Newey, and Robins (2018), ``Double/debiased machine learning for treatment and
structural parameters'', \emph{Econometrics Journal}}{The DML framework, orthogonality, and cross-fitting}{Sections 1--3}{3}
\reading{Chernozhukov, Hansen, Kallus, Spindler, and Syrgkanis, \emph{Applied Causal Inference Powered by ML and AI} (free online)}{DML at
textbook pace, with code}{Ch.\ 4, 9--10}{2}
\reading{Cinelli and Hazlett (2020), ``Making sense of sensitivity'', \emph{JRSS B}}{Partial-$R^2$ sensitivity, robustness values, and
benchmarking}{Sections 1--4}{2}
\reading{Chernozhukov, Cinelli, Newey, Sharma, and Syrgkanis, ``Long story short: omitted variable bias in causal machine learning'' (working
paper)}{The sensitivity bound for DML and the interactive model}{Sections 1--3}{3}
\reading{Belloni, Chernozhukov, and Hansen (2014), ``Inference on treatment effects after selection among high-dimensional controls'',
\emph{Review of Economic Studies}}{Post-double-selection lasso}{Sections 1--2}{3}
\reading{Handout H5, \emph{Regression discontinuity}}{Identification from a threshold rule}{All}{2}
\end{readinglist}

\end{document}
